\documentclass[12pt]{article}
\usepackage[T1]{fontenc}
\usepackage[utf8]{inputenc}
\usepackage{lmodern}
\usepackage{textcomp}
\usepackage{enumitem}
\usepackage{geometry}
\usepackage{graphicx}
\usepackage{amsmath, amssymb}
\geometry{margin=1in}
\begin{document}
\begin{center}
Constitutional Law - Undergraduate Level Assessment 16 \
Total Marks: 40 \
Answer all questions.
\end{center}

\section*{Multiple Choice (10 marks)}
\begin{enumerate}[label=\arabic*.]
\item Which statement is closest to characterizing the principal difference between the positions adopted by Hobbes and Locke?
\begin{enumerate}[label=(\alph*)]
    \item They adopt different attitudes towards the role of the courts in maintaining order.
    \item They disagree about the role of law in society.
    \item They have opposing views about the nature of contractual obligations.
    \item They differ in respect of their account of life before the social contract.
\end{enumerate}
\item How does Nozick answer the criticism of his historical entitlement theory that if the distribution of goods in society is unjust those at the bottom always lose?
\begin{enumerate}[label=(\alph*)]
    \item It can be remedied by redistribution of wealth.
    \item If each person's holdings are just, then the total distribution of holdings is just.
    \item Historical factors are secondary to moral imperatives.
    \item He has no answer.
\end{enumerate}
\item Which of the following most accurately describes Hart's response to Fuller's argument concerning the invalidity of Nazi law?
\begin{enumerate}[label=(\alph*)]
    \item The Nazi law in question was validly enacted.
    \item The court misunderstood the legislation.
    \item Fuller misconstrued the purpose of the law.
    \item The Nazi rule of recognition was unclear.
\end{enumerate}
\item What is passive personality jurisdiction?
\begin{enumerate}[label=(\alph*)]
    \item It is jurisdiction based on the nationality of the offender
    \item It is jurisdiction based on where the offence was committed
    \item It is jurisdiction based on the nationality of the victims
    \item It is jurisdiction based on the country where the legal person was Registered
\end{enumerate}
\item Who said that "Jurisprudence is the eye of law"
\begin{enumerate}[label=(\alph*)]
    \item Maine
    \item Savigny
    \item Pound
    \item Laski
\end{enumerate}
\end{enumerate}

\section*{True / False (10 marks)}
\textit{Answer True or False}
\begin{enumerate}[label=\arabic*.]
\setcounter{enumi}{5}
\item True or False: Jurisdiction is ordinarily determined by the location of the offender, regardless of where the offense occurred.
\item True or False: John Austin wrote 'The Concept of Law', which outlines his legal positivism theory.
\item True or False: The people in the original position would likely choose equality above liberty as a foundational principle.
\item True or False: Austin argued that law is merely a collection of commands issued by a sovereign authority without regard to justice.
\item True or False: Third generation rights are considered more important than civil and political rights according to the UN Vienna Declaration.
\end{enumerate}

\section*{Long Form Response (20 marks)}
\textit{Answer each question with a clear and complete explanation.}
\begin{enumerate}[label=\arabic*.]
\setcounter{enumi}{10}
\item Discuss the characterization of Austin as a 'naive empiricist' in constitutional law, focusing on his pragmatic approach to understanding laws rather than a conceptual one.
\item Discuss Ronald Dworkin's criticism of Richard Posner's economic analysis of law, particularly his assertion that wealth cannot be regarded as a fundamental value. What are the implications of this critique?
\end{enumerate}

\vfill
\noindent\textit{End of Paper}
\end{document}